\documentclass[11pt,a4paper]{scrartcl}

\usepackage{amssymb}
\usepackage{amsmath}

\usepackage[utf8]{inputenc}
\usepackage[T1]{fontenc}
\newcommand{\ats}[1]{\par\noindent\textit{Atsakymas:} #1}
\usepackage{cancel}
\usepackage{tikz}
\usepackage{lmodern} 
\usepackage{amsmath,amssymb,amsthm}
\usepackage{graphicx}
\usepackage[dvipsnames,svgnames]{xcolor}
\usepackage{hyperref}
\usepackage{mathrsfs}
\usepackage[shortlabels]{enumitem}
\usepackage[obeyFinal,textsize=scriptsize,shadow,loadshadowlibrary]{todonotes}
\usepackage{mathtools}
\usepackage{microtype}

%%% Paragraph layout
\setlength{\parindent}{0pt}
\setlength{\parskip}{0.6\baselineskip}

%%% Colored sections
\renewcommand*{\sectionformat}%
  {\color{purple}\S\thesection\enskip}
\renewcommand*{\subsectionformat}%
  {\color{purple}\S\thesubsection\enskip}
\renewcommand*{\subsubsectionformat}%
  {\color{purple}\S\thesubsubsection\enskip}
\addtokomafont{paragraph}{\color{orange!35!black}\P\ }
\KOMAoptions{numbers=noenddot}

%%% Title
\addtokomafont{subtitle}{\Large}
\setkomafont{author}{\Large\scshape}
\setkomafont{date}{\Large\normalsize}
% Neigiamas vertikalus tarpas, kuris pakelia dokumento antraštę į viršų. 
% Galite keisti -2.5cm į -3cm (aukščiau) arba -1.5cm (žemiau).
\titlehead{\vspace*{-7.5cm}} 

% PRIDĖTA: Automatiškai perrašome datą į mokyklos pavadinimą
\AtBeginDocument{%
  \date{\normalsize Vilniaus licėjus}
}

%%% Page setup
\usepackage[headsepline]{scrlayer-scrpage}
\renewcommand{\headfont}{}
\addtolength{\textheight}{3.14cm}
\setlength{\footskip}{0.5in}
\setlength{\headsep}{10pt}
% Puslapio antraštė turi rodyti tik konkrečius dokumento duomenis.
% Nustatome juos aiškiai, kad neatsirastų "\@author" / "\@title" arba senų reikšmių.
\makeatletter
\AtBeginDocument{%
  \ihead{\footnotesize\textbf{Andrius Berniukevičius}}%
  \ohead{\footnotesize\textbf{\@title}}%
}
\makeatother
\automark{section}
\chead{}
\cfoot{\pagemark}

%%% Macros
\providecommand{\ol}{\overline}
\providecommand{\ul}{\underline}
\providecommand{\wt}{\widetilde}
\providecommand{\wh}{\widehat}
\providecommand{\eps}{\varepsilon}
\providecommand{\half}{\frac{1}{2}}
\providecommand{\inv}{^{-1}}
\newcommand{\dang}{\measuredangle} %% Directed angle
\newcommand{\vocab}[1]{\textbf{\color{blue}\sffamily #1}}
\providecommand{\CC}{\mathbb C}
\providecommand{\FF}{\mathbb F}
\providecommand{\NN}{\mathbb N}
\providecommand{\QQ}{\mathbb Q}
\providecommand{\RR}{\mathbb R}
\providecommand{\ZZ}{\mathbb Z}
\providecommand{\ts}{\textsuperscript}
\providecommand{\dg}{^\circ}
\providecommand{\ii}{\item}
\providecommand{\defeq}{\coloneqq}
\DeclareMathOperator*{\lcm}{lcm}
\DeclareMathOperator*{\argmin}{arg min}
\DeclareMathOperator*{\argmax}{arg max}
\providecommand{\hrulebar}{\par
\hspace{\fill}\rule{0.95\linewidth}{.7pt}\hspace{\fill}
\par\nointerlineskip \vspace{\baselineskip}}

%%% Optional colored theorems
\usepackage{thmtools}
\usepackage[framemethod=TikZ]{mdframed}
\mdfdefinestyle{mdbluebox}{%
  roundcorner=10pt,
  linewidth=1pt,
  skipabove=12pt,
  innerbottommargin=9pt,
  skipbelow=2pt,
  linecolor=blue,
  nobreak=true,
  backgroundcolor=TealBlue!5,
}
\declaretheoremstyle[
  headfont=\sffamily\bfseries\color{MidnightBlue},
  mdframed={style=mdbluebox},
  headpunct={\\[3pt]},
  postheadspace={0pt}
]{thmbluebox}
\mdfdefinestyle{mdredbox}{%
  linewidth=0.5pt,
  skipabove=12pt,
  frametitleaboveskip=5pt,
  frametitlebelowskip=0pt,
  skipbelow=2pt,
  frametitlefont=\bfseries,
  innertopmargin=4pt,
  innerbottommargin=8pt,
  nobreak=true,
  backgroundcolor=Salmon!5,
  linecolor=RawSienna,
}
\declaretheoremstyle[
  headfont=\bfseries\color{RawSienna},
  mdframed={style=mdredbox},
  headpunct={\\[3pt]},
  postheadspace={0pt},
]{thmredbox}
\mdfdefinestyle{mdgreenbox}{%
  skipabove=8pt,
  linewidth=2pt,
  rightline=false,
  leftline=true,
  topline=false,
  bottomline=false,
  linecolor=ForestGreen,
  backgroundcolor=ForestGreen!5,
}
\declaretheoremstyle[
  headfont=\bfseries\sffamily\color{ForestGreen!70!black},
  bodyfont=\normalfont,
  spaceabove=2pt,
  spacebelow=1pt,
  mdframed={style=mdgreenbox},
  headpunct={ --- },
]{thmgreenbox}
\mdfdefinestyle{mdblackbox}{%
  skipabove=8pt,
  linewidth=3pt,
  rightline=false,
  leftline=true,
  topline=false,
  bottomline=false,
  linecolor=black,
  backgroundcolor=RedViolet!5!gray!5,
}
\declaretheoremstyle[
  headfont=\bfseries,
  bodyfont=\normalfont\small,
  spaceabove=0pt,
  spacebelow=0pt,
  mdframed={style=mdblackbox}
]{thmblackbox}

%% Aplinkos su rėmeliais (thmtools)
\declaretheorem[style=thmbluebox,name=Teorema]{theorem}
\declaretheorem[style=thmbluebox,name=Lema]{lemma}
\declaretheorem[style=thmbluebox,name=Savybė]{property}
\declaretheorem[style=thmbluebox,name=Teiginys]{proposition}
\declaretheorem[style=thmbluebox,name=Išvada]{corollary}
\declaretheorem[style=thmbluebox,name=Teorema,numbered=no]{theorem*}
\declaretheorem[style=thmbluebox,name=Lema,numbered=no]{lemma*}
\declaretheorem[style=thmbluebox,name=Teiginys,numbered=no]{proposition*}
\declaretheorem[style=thmbluebox,name=Išvada,numbered=no]{corollary*}
\declaretheorem[style=thmgreenbox,name=Algoritmas]{algorithm}
\declaretheorem[style=thmgreenbox,name=Algoritmas,numbered=no]{algorithm*}
\declaretheorem[style=thmgreenbox,name=Teiginys]{claim}
\declaretheorem[style=thmgreenbox,name=Teiginys,numbered=no]{claim*}
\declaretheorem[style=thmredbox,name=Pavyzdys]{example}
\declaretheorem[style=thmredbox,name=Pavyzdys,numbered=no]{example*}
\declaretheorem[style=thmblackbox,name=Pastaba]{remark}
\declaretheorem[style=thmblackbox,name=Pastaba,numbered=no]{remark*}

%% Standartinės amsthm aplinkos (Apibrėžimai ir užduotys)
\theoremstyle{definition}
\ifcsname conjecture\endcsname\else\newtheorem{conjecture}{Hipotezė}\fi
\ifcsname definition\endcsname\else\newtheorem{definition}{Apibrėžimas}\fi
\ifcsname fact\endcsname\else\newtheorem{fact}{Faktas}\fi
\ifcsname answer\endcsname\else\newtheorem{answer}{Atsakymas}\fi
\ifcsname ques\endcsname\else\newtheorem{ques}{Klausimas}\fi
\ifcsname exercise\endcsname\else\newtheorem{exercise}{Užduotis}\fi
\ifcsname problem\endcsname\else\newtheorem{problem}{Uždavinys}[section]\fi
\ifcsname conjecture*\endcsname\else\newtheorem*{conjecture*}{Hipotezė}\fi
\ifcsname definition*\endcsname\else\newtheorem*{definition*}{Apibrėžimas}\fi
\ifcsname fact*\endcsname\else\newtheorem*{fact*}{Faktas}\fi
\ifcsname answer*\endcsname\else\newtheorem*{answer*}{Atsakymas}\fi
\ifcsname ques*\endcsname\else\newtheorem*{ques*}{Klausimas}\fi
\ifcsname exercise*\endcsname\else\newtheorem*{exercise*}{Užduotis}\fi
\ifcsname problem*\endcsname\else\newtheorem*{problem*}{Uždavinys}\fi

\newcounter{answerssection}
\newcommand{\answerssection}{%
  \setcounter{answerssection}{\value{section}}%
  \section{Atsakymai}%
  \setcounter{problem}{0}%
  \renewcommand{\theproblem}{\arabic{answerssection}.\arabic{problem}}%
}

%% GALUTINIS SPRENDIMAS PROOF APLINKAI
%% --- PRADŽIA: Įrodymo aplinkos sutvarkymas ---
\makeatletter

% 1. Užtikriname, kad žodis būtų "Įrodymas", net jei "babel" paketas bando jį perrašyti
\AtBeginDocument{%
  \renewcommand{\proofname}{Įrodymas}%
  \@ifpackageloaded{babel}{%
    \addto\captionslithuanian{\renewcommand{\proofname}{Įrodymas}}%
  }{}%
}

% 2. Priverstinai pašaliname scrartcl klasės bold šriftą ir paliekame TIK italic
% 3. Sujungiame su tuščiu kvadratėliu dešiniajame kampe.
\renewcommand{\qedsymbol}{\ensuremath{\Box}}
\renewenvironment{proof}[1][\proofname]{\par
  \pushQED{\qed}%
  \normalfont \topsep6\p@\@plus6\p@\relax
  \trivlist
  \item[\hskip\labelsep
        \normalfont\itshape % Priverčia naudoti tik pasvirą šriftą, kaip nuotraukoje
    #1\@addpunct{.}]\ignorespaces
}{%
  \popQED\endtrivlist\@endpefalse
}
\newenvironment{solution}{\par
  \normalfont \topsep6\p@\@plus6\p@\relax
  \trivlist
  \item[\hskip\labelsep
        \normalfont\itshape
    \textit{Sprendimas}\@addpunct{.}]\ignorespaces
}{%
  \endtrivlist\@endpefalse
}
\newcommand{\ans}[1]{\par\noindent\textit{Atsakymas:} #1}
\makeatother


\title{Algebra II: Vijeto teorema ir simetrija}
\author{Andrius Berniukevičius}

\newenvironment{solution}
    {\begin{list}{}{\setlength{\leftmargin}{10pt}\setlength{\rightmargin}{10pt}}\item[]}
    {\end{list}}

\begin{document}

\maketitle

\section{Kaip apgauti Diskriminantą?}

Mokykloje jūs išmokote galingą įrankį – diskriminantą ($D = b^2 - 4ac$). Su juo galite rasti bet kokios kvadratinės lygties $ax^2 + bx + c = 0$ šaknis. 

Tačiau įsivaizduokite olimpiadinį uždavinį: duota lygtis $x^2 - 6x + 2 = 0$. Prašoma rasti jos šaknų kvadratų sumą ($x_1^2 + x_2^2$). 
Jei skaičiuosite diskriminantą, gausite $D = 36 - 8 = 28$. Šaknys bus negražios: $x_1 = \frac{6 - \sqrt{28}}{2}$ ir $x_2 = \frac{6 + \sqrt{28}}{2}$. Kelti tokius monstrus kvadratu ir juos sudėti yra baisus darbas, kuriame labai lengva padaryti klaidą.

Olimpiadinėje matematikoje mes šaknų \textbf{neieškome}. Mes ieškome būdo, kaip sužinoti atsakymą išvengiant juodo darbo. Tam mums padeda prancūzų matematiko Fransua Vijeto (François Viète) atradimas.

\textbf{Vijeto teorema:}
Jei kvadratinės lygties pavidalo $\mathbf{x^2 + px + q = 0}$ šaknys yra $x_1$ ir $x_2$, tai:
\begin{itemize}
    \item Šaknų \textbf{suma} visada lygi koeficientui prie $x$ su priešingu ženklu: $\mathbf{x_1 + x_2 = -p}$
    \item Šaknų \textbf{sandauga} visada lygi laisvajam nariui: $\mathbf{x_1 \cdot x_2 = q}$
\end{itemize}
\textit{Pastaba: Jei lygtis prasideda skaičiumi, pvz., $ax^2 + bx + c = 0$, pirmiausia visą lygtį padaliname iš $a$, kad prie $x^2$ liktų vienetas: $x^2 + \frac{b}{a}x + \frac{c}{a} = 0$. Tuomet $x_1 + x_2 = -\frac{b}{a}$, o $x_1 x_2 = \frac{c}{a}$.}

% ==========================================
% SIMETRINIAI REIŠKINIAI
% ==========================================
\section{Algebrinė simetrija}

Vijeto teorema mums duoda tik šaknų sumą ir sandaugą. Bet ką daryti, jei uždavinys prašo rasti $x_1^2 + x_2^2$ arba $\frac{1}{x_1} + \frac{1}{x_2}$?

Čia į pagalbą ateina greitosios daugybos formulės. Kadangi reiškiniai yra \textbf{simetriniai} (sukeitus $x_1$ ir $x_2$ vietomis, reiškinys nepasikeičia), mes juos visada galime pertvarkyti taip, kad liktų \textbf{tik} sumos ir sandaugos.

\textbf{1. Kvadratų suma: $x_1^2 + x_2^2$}
Iš praeito užsiėmimo atsimename pilno kvadrato formulę: $(x_1 + x_2)^2 = x_1^2 + 2x_1x_2 + x_2^2$.
Iš čia labai lengva išreikšti plikų kvadratų sumą, tiesiog atėmus vidurinį narį!
\[ \mathbf{x_1^2 + x_2^2 = (x_1 + x_2)^2 - 2x_1x_2} \]
Matote? Reiškinys dabar sudarytas tik iš Vijeto sumos ir Vijeto sandaugos!

\textbf{2. Trupmenų suma: $\frac{1}{x_1} + \frac{1}{x_2}$}
Subendravardikliname (pirmą trupmeną dauginame iš $x_2$, antrą iš $x_1$):
\[ \mathbf{\frac{1}{x_1} + \frac{1}{x_2} = \frac{x_1 + x_2}{x_1 x_2}} \]
Vėlgi, viršuje Vijeto suma, apačioje – sandauga.

\textbf{3. Kubų suma: $x_1^3 + x_2^3$}
Pritaikome kubų sumos formulę (iš pirmojo handouto) ir tuomet į kvadratų sumą įstatome triuką su minus $2x_1x_2$:
\[ x_1^3 + x_2^3 = (x_1 + x_2)(x_1^2 - x_1x_2 + x_2^2) = \mathbf{(x_1 + x_2)\big((x_1 + x_2)^2 - 3x_1x_2\big)} \]

% ==========================================
% ŠAKNIES ĮSTATYMO TRIUKAS
% ==========================================
\section{Šaknies įstatymo triukas}

Kartais olimpiadose prašoma apskaičiuoti reiškinį, kuris nėra simetriškas (pvz., $x_1^2 - 5x_1$). Tokiu atveju veikia pats primityviausias, bet genialus triukas, kurį mokiniai dažnai pamiršta: \textbf{šaknis privalo tenkinti lygtį}.

Jei $x_1$ yra lygties $x^2 - 6x + 2 = 0$ šaknis, vadinasi, mes drąsiai galime į lygtį vietoj $x$ įstatyti $x_1$:
\[ x_1^2 - 6x_1 + 2 = 0 \]
Iš šios lygybės galime išreikšti $x_1^2$:
\[ \mathbf{x_1^2 = 6x_1 - 2} \]
Tai leidžia mums bet kokį kvadratą „nuleisti“ į pirmo laipsnio reiškinį!

% ==========================================
% PAVYZDŽIAI
% ==========================================
\section{Olimpiadiniai pavyzdžiai}

\begin{example}[Sistemos apgavystė]
Lygties $x^2 - 6x + 2 = 0$ šaknys yra $x_1$ ir $x_2$. Nespręsdami lygties, raskite reiškinio $x_1^2 + x_2^2$ reikšmę.
\end{example}

\textbf{Sprendimas:}\\
Iš Vijeto teoremos žinome, kad:
$x_1 + x_2 = 6$ (koeficientas prie $x$ su priešingu ženklu)
$x_1 \cdot x_2 = 2$ (laisvasis narys)

Pertvarkome prašomą rasti reiškinį naudodami pilno kvadrato formulę:
\[ x_1^2 + x_2^2 = (x_1 + x_2)^2 - 2x_1x_2 \]
Įstatome mūsų Vijeto reikšmes:
\[ x_1^2 + x_2^2 = 6^2 - 2 \cdot 2 = 36 - 4 = \mathbf{32} \hfill \square \]

\begin{example}[Naujos lygties sudarymas]
Sukurkite (parašykite) naują kvadratinę lygtį, kurios šaknys būtų atvirkštinės lygties $x^2 - 5x + 3 = 0$ šaknims (t.y. $\frac{1}{x_1}$ ir $\frac{1}{x_2}$).
\end{example}

\textbf{Sprendimas:}\\
Naujoji lygtis atrodys taip: $x^2 + px + q = 0$. Jos koeficientai $p$ ir $q$ priklausys nuo naujųjų šaknų:
Pagal Vijetą: $-p = \frac{1}{x_1} + \frac{1}{x_2}$ ir $q = \frac{1}{x_1} \cdot \frac{1}{x_2}$.

Iš pradinės lygties žinome, kad $x_1 + x_2 = 5$ ir $x_1x_2 = 3$.
Apskaičiuokime naująjį $q$:
\[ q = \frac{1}{x_1} \cdot \frac{1}{x_2} = \frac{1}{x_1x_2} = \frac{1}{3} \]
Apskaičiuokime naująjį $-p$:
\[ -p = \frac{1}{x_1} + \frac{1}{x_2} = \frac{x_1 + x_2}{x_1x_2} = \frac{5}{3} \implies p = -\frac{5}{3} \]
Mūsų naujoji lygtis yra: $x^2 - \frac{5}{3}x + \frac{1}{3} = 0$.
Kad būtų gražiau, padauginkime viską iš 3:
\[ \mathbf{3x^2 - 5x + 1 = 0} \hfill \square \]
\textit{Olimpiadinis pastebėjimas: Jei norime lygties su atvirkštinėmis šaknimis, užtenka tiesiog sukeisti vietomis koeficientus $a$ ir $c$!}

\begin{example}[Asimetrinis reiškinys]
Lygties $x^2 - 5x - 2 = 0$ šaknys yra $x_1$ ir $x_2$. Raskite reiškinio $x_1^2 - 4x_1 + x_2$ reikšmę.
\end{example}
\textbf{Sprendimas:}\\
Reiškinys nesimetriškas. Negalime iškart taikyti Vijeto, nes $x_1$ pakeltas kvadratu, o $x_2$ – ne. Taikome \textbf{Šaknies įstatymo triuką}.
Kadangi $x_1$ yra šios lygties šaknis, jis paverčia ją teisinga lygybe:
\[ x_1^2 - 5x_1 - 2 = 0 \]
Iš čia išsireiškiame $x_1^2$:
\[ x_1^2 = 5x_1 + 2 \]
Dabar šią išraišką įstatome į ieškomą reiškinį vietoj $x_1^2$:
\[ (5x_1 + 2) - 4x_1 + x_2 \]
Sutraukiame panašius narius ($5x_1 - 4x_1$):
\[ x_1 + x_2 + 2 \]
Magija! Liko tik šaknų suma, kurią mes žinome iš Vijeto teoremos ($x_1 + x_2 = 5$).
Įstatome ir gauname atsakymą:
\[ 5 + 2 = \mathbf{7} \hfill \square \]

\vspace{0.5cm}

% ==========================================
% UŽDAVINIAI
% ==========================================
\newpage
\section{Uždaviniai savarankiškam sprendimui (30 iššūkių)}

\begin{enumerate}
    \renewcommand{\labelenumi}{\textbf{\theenumi.}}
    
    \item Lygties $x^2 - 7x + 10 = 0$ šaknys yra $x_1$ ir $x_2$. Nespręsdami lygties, užrašykite šaknų sumą ir šaknų sandaugą.
    \item Lygties $2x^2 + 8x - 5 = 0$ šaknys yra $x_1$ ir $x_2$. Raskite $x_1 + x_2$ ir $x_1 \cdot x_2$.
    \item Lygties $3x^2 - 12x = 0$ šaknys yra $x_1$ ir $x_2$. Nespręsdami lygties (nors čia ir lengva), užrašykite Vijeto sumą ir sandaugą.
    \item Lygties $x^2 + 5x - 8 = 0$ šaknys yra $x_1$ ir $x_2$. Nespręsdami lygties apskaičiuokite reiškinio $(x_1 + 1)(x_2 + 1)$ reikšmę.
    \item Lygties $x^2 - 9x + 4 = 0$ šaknys yra $x_1$ ir $x_2$. Nespręsdami lygties, raskite $x_1^2 + x_2^2$.
    \item Lygties $x^2 + 4x - 6 = 0$ šaknys yra $x_1$ ir $x_2$. Raskite reiškinio $\frac{1}{x_1} + \frac{1}{x_2}$ reikšmę.
    \item Lygties $2x^2 - 5x + 1 = 0$ šaknys yra $x_1$ ir $x_2$. Raskite reiškinio $\frac{3}{x_1} + \frac{3}{x_2}$ reikšmę.
    
    \item Lygties $x^2 - 5x + 2 = 0$ šaknys yra $x_1$ ir $x_2$. Nespręsdami lygties, raskite $(x_1 - x_2)^2$.
    \item Naudodamiesi praeito uždavinio duomenimis (lygčiai $x^2 - 5x + 2 = 0$), raskite $\frac{x_1}{x_2} + \frac{x_2}{x_1}$.
    \item Naudodamiesi tos pačios lygties duomenimis, raskite $x_1^3 + x_2^3$.
    \item Lygties $x^2 - 3x - 2 = 0$ šaknys yra $x_1$ ir $x_2$. Raskite reiškinio $x_1^2 x_2 + x_1 x_2^2$ reikšmę.
    \item Lygties $x^2 - 4x + 1 = 0$ šaknys yra $x_1$ ir $x_2$. Raskite reiškinio $\frac{x_1^2}{x_2} + \frac{x_2^2}{x_1}$ reikšmę.
    \item Sukurkite (parašykite) kvadratinę lygtį, kurios šaknys būtų $3 - \sqrt{5}$ ir $3 + \sqrt{5}$.
    \item Lygties $x^2 - 5x + 3 = 0$ šaknys yra $x_1$ ir $x_2$. Sukurkite naują kvadratinę lygtį, kurios šaknys būtų $(x_1 + 2)$ ir $(x_2 + 2)$.
    \item Ten pat. Sukurkite naują kvadratinę lygtį, kurios šaknys būtų atvirkštiniai skaičiai pradinėms šaknims: $\frac{1}{x_1}$ ir $\frac{1}{x_2}$.
    \item Lygties $x^2 - 4x + 1 = 0$ šaknys yra $x_1$ ir $x_2$. Sukurkite naują lygtį, kurios šaknys būtų $x_1^2$ ir $x_2^2$.
    
    \item Su kokia parametro $a$ reikšme lygties $x^2 + (a-2)x + (a+5) = 0$ šaknų suma yra lygi lygties šaknų sandaugai?
    \item Su kokia parametro $k$ reikšme lygties $x^2 - 6x + k = 0$ viena šaknis yra dvigubai didesnė už kitą ($x_1 = 2x_2$)?
    \item Raskite visas parametro $m$ reikšmes, su kuriomis lygties $x^2 - mx + 12 = 0$ šaknų skirtumas yra lygus lygiai $1$.
    \item Raskite parametrą $p$, jeigu žinoma, kad lygties $x^2 + px + 8 = 0$ šaknų kvadratų suma lygi $20$.
    \item Kvadratinės lygties $ax^2 + bx + c = 0$ viena šaknis yra kitos šaknies atvirkštinis skaičius (t.y. $x_2 = \frac{1}{x_1}$). Kokią sąlygą privalo tenkinti šios lygties koeficientai?
    \item Lygties $x^2 - 14x + 9 = 0$ šaknys yra teigiamos. Raskite reiškinio $\sqrt{x_1} + \sqrt{x_2}$ reikšmę.
    \item Lygties $x^2 - 8x + 4 = 0$ šaknys yra teigiamos. Raskite reiškinio $|\sqrt{x_1} - \sqrt{x_2}|$ reikšmę.
    \item Raskite visas parametro $a$ reikšmes, su kuriomis lygties $x^2 - ax + 3 = 0$ šaknų kvadratų suma būtų lygi $10$.
    
    \item Lygties $x^2 - 3x - 5 = 0$ šaknys yra $x_1$ ir $x_2$. Raskite reiškinio $x_1^2 - 2x_1 + x_2$ reikšmę.
    \item Lygties $x^2 - 7x + 3 = 0$ šaknys yra $x_1$ ir $x_2$. Raskite reiškinio $x_1^3 - 7x_1^2 + 5x_1 + 2x_2$ reikšmę.
    \item Lygties $x^2 - x - 3 = 0$ šaknys yra $x_1$ ir $x_2$. Apskaičiuokite reiškinio $x_1^3 + 4x_2^2 + 10$ reikšmę.
    \item Tegu $x_1$ ir $x_2$ yra lygties $x^2 + x - 1 = 0$ šaknys. Apskaičiuokite $x_1^8 + x_2^8$.
    \item Raskite visas parametro $m$ reikšmes, su kuriomis lygties $x^2 - mx + m - 1 = 0$ šaknų kvadratų suma yra pati mažiausia galima. Kokia tai suma?
    \item Yra žinoma, kad lygtis $x^2 + px + q = 0$ turi dvi šaknis $x_1$ ir $x_2$, o lygtis $x^2 - px - q = 0$ (koeficientų ženklai sukeisti į minusus) turi šaknis $x_3$ ir $x_4$. Įrodykite, kad šaknų suma $x_1 + x_2 + x_3 + x_4 = 0$, ir suraskite reiškinio $x_1^2 + x_2^2 - x_3^2 - x_4^2$ reikšmę (atsakymas priklausys nuo $p$ ir $q$).
\end{enumerate}

\end{document}